\section{Implementation}
\label{sec:implementation}

\subsection{Two method instances and their controls}
We implement the construction in modified llama.cpp paths~\cite{records}. The phase executor owns a resident target context with a dense, contiguous, nonwrapping sequence and F16 keys and values. A backend-neutral planner preserves ordered rows, physical extents and requested-output flags; native adapters supply numerical descriptors and operations. The CUDA instance uses causal, non-Flash attention for the supported Llama and Qwen architectures.

For Q8 projections, full arithmetic groups have $a=4$. The packed path retains the existing small-column arithmetic~\cite{loop4code}, presents complete groups as channels, and restores column order afterward. The \emph{backend-loop4} control already quantizes the wide activation once before invoking four-column and tail kernels. The comparison therefore tests execution organization beyond shared quantization. Cross-group packing requires at least eight known rows. In the $k=3$ speculative applications, target verification has four rows, so this packing opportunity occurs in wider prefill work. The \Simple/\Local{} factor separately controls placement of the physical attention cuts.

The 32B mixed Q4\_K/Q6\_K instance instead uses $a=1$ and whole-graph cuts. B shares activation quantization and loops over scalar backend projections; C presents a $[K,1,q]$ view and reshapes the result. Both retain the one-column dot-product specialization. The S3 C/B comparison consequently measures scheduling of independent scalar groups, distinct from the Q8 \Local-plus-packed combination. The weaker G control performs scalar slicing at graph level.

The Vulkan instance retains F32 activations where the scalar reference uses F32, together with the required attention extent and layout. It provides a separate implementation of the grouping principle, with its own numerical reference and controlled append measurements.

\subsection{Output ownership, reuse and implementation coverage}
After each native call, the executor synchronizes and copies requested vocabulary vectors into storage indexed by logical row. It then checks the physical append endpoint. This preserves an earlier accepted row when subsequent calls reuse native output buffers. Planning uses native metadata without reference replay or KV export.

Graph reuse compares the numerical policy, complete per-row extent vector, relevant head-normalization and rotary-placement choices, and the engine's topology and output-count conditions. Requested row indices are refreshed as inputs. A mismatch rebuilds the graph. The deployment fixes additional numerical modules and model choices; the frontend's cache-type tag alone does not encode them all.

Online admission checks the supported cache and output reserve. Marked backend-loop4 operators additionally check descriptors, actual device residency and dispatch constraints. Packed nodes use ordinary backend dispatch and do not inherit every marker check; their arithmetic evidence comes from deployment-specific qualification. Capture runs collect selected operators, full vectors and native state for offline comparison, while clean service timings exclude those observers. The supplement maps these mechanisms to the actual source paths.

A failed target group poisons its context. The published command-line worker contains other failures by terminating, rather than recovering a partially updated service transaction. Detailed failure coverage is in the supplement. The method's state-commit argument concerns successful admitted executions; a reusable service library would also need transaction-wide invalidation on failure.
